La brossa espacial deixada per satèl·lits antics i els seus coets llançadors s'està convertint en un perill per a altres satèl·lits. El novembre de 2023, durant unes tasques de reparació, dos astronautes es van deixar una caixa d'eines a l'exterior de l'\textit{Estació Espacial Internacional} (EEI).

\begin{apartat}{1,25}
A partir de la llei de gravitació universal, deduïu l'expressió de la velocitat orbital en funció del radi orbital. Calculeu la velocitat de la caixa d'eines en òrbita a 400~km per sobre de la superfície terrestre i el nombre de voltes que farà la caixa cada dia al voltant de la Terra.
\end{apartat}

\begin{apartat}{1,25}
La fi de l'EEI està planificada per a l'any 2031. D'una manera gradual i controlada se'n baixarà l'òrbita fins als 280~km d'altura per sobre de la superfície terrestre. Calculeu l'energia mecànica de l'EEI en aquesta òrbita i justifiqueu-ne el signe. L'última tripulació abandonarà l'estació i, posteriorment, l'estació caurà des d'aquesta altura al mig de l'oceà Pacífic. Calculeu amb quina energia cinètica impactarà l'estació contra l'aigua, sense tenir en compte els efectes de l'atmosfera terrestre.
\end{apartat}

\dades{%
  $G = 6,67\times 10^{-11}\ \mathrm{N\,m^2\,kg^{-2}}$.\\
  Massa de la Terra: $M_\mathrm{T} = 5,98\times 10^{24}\ \mathrm{kg}$.\\
  Radi de la Terra: $R_\mathrm{T} = 6,37\times 10^{6}\ \mathrm{m}$.\\
  Massa de l'EEI: $M_\mathrm{EEI} = 430\times 10^{3}\ \mathrm{kg}$.}
